\section{Entrenamiento a gran escala: fallas numéricas, factores de confusión y el iceberg de I\&D}

Cuando el entrenamiento llega a la escala de miles de GPU y semanas de ejecución, el éxito de los experimentos pequeños ya no basta para demostrar que una configuración es segura. En clase se dio un ejemplo: cierta ruta en FP16 parecía coincidir con FP32 en modelos más pequeños, pero en el entrenamiento de un modelo grande la loss empezó a subir después de unas dos semanas. Lo más difícil no es «saber que FP16 puede ser inestable», sino que, cuando ocurre la falla, coexisten numerosos \term{confounding factors} (factores de confusión): datos, optimizador, red, hardware, código y aleatoriedad.

\subsection{Por qué la validación a pequeña escala no detecta errores de gran escala}

\term{FP16} es punto flotante de 16 bits (16-bit floating point); ahorra memoria de video y ancho de banda, pero el rango del exponente y de la mantisa es limitado. El error puede ser pequeño en una sola capa, en secuencias cortas o en entrenamientos breves, y aun así amplificarse gradualmente con redes más profundas, un horizon más largo y una mayor varianza del gradiente. La práctica segura no consiste en volver por completo a FP32, sino en identificar los operadores sensibles y usar BF16, FP32 accumulation, loss scaling, monitoreo numérico y validación de escala por etapas.

\lecturefigure{06-scale-amplifies-bugs.png}{Una aproximación numérica que pasa a pequeña escala puede convertirse en una falla diferida en un entrenamiento grande y prolongado.}{Subtítulos locales 12:08--13:35; redibujo conceptual.}

\begin{knowledgebox}{Lectura de la figura: la exposición diferida cambia el diseño experimental}
Si un error aparece hasta dos semanas después, una sola búsqueda binaria para localizarlo puede consumir miles de GPU-day. El equipo necesita registrar desde las primeras etapas del entrenamiento las normas del gradiente, los rangos de activación, la loss desagregada por grupos, los errores de hardware, los shards de datos y las versiones del código, y diseñar un canary de escala media que amplifique los riesgos lo antes posible. La observabilidad no es un informe posterior al entrenamiento, sino un instrumento experimental que reduce el espacio de búsqueda de la depuración.
\end{knowledgebox}

\subsection{Por qué el final run cost publicado es engañoso}

Los artículos suelen reportar el número de GPU, los tokens y los días del entrenamiento final; esas cifras solo describen la ejecución principal exitosa y no incluyen las pruebas de receta, los checkpoint fallidos, la limpieza de datos, la optimización de kernels ni la construcción del sistema. La discusión de costos de LLaMA y DeepSeek en clase advierte que multiplicar el cómputo publicado por el precio de la nube solo da una factura aproximada de una ejecución; no permite deducir el costo total de I\&D que le llevó al equipo obtener esa receta.

\lecturefigure{07-rd-iceberg.png}{El entrenamiento final exitoso es solo el costo visible; lo verdaderamente caro es todo el proceso para obtener una receta estable.}{Subtítulos locales 14:20--17:05; las estimaciones de costo del ponente son solo contexto de la clase.}

\begin{warningbox}{No tome las estimaciones del ponente como datos de auditoría}
El número de GPU, el precio por hora y el costo total de seis meses mencionados en la entrevista se dijeron de viva voz, y los subtítulos podrían confundir A100 con H100. La nota conserva la conclusión didáctica «la escala es muy grande, el ensayo y error es muy caro y el final run subestima la I\&D», y no usa estas estimaciones para extraer conclusiones financieras precisas.
\end{warningbox}

\subsection{Memoria organizacional: por qué importa «haberlo hecho una vez»}

El know-how del entrenamiento a gran escala incluye mucho conocimiento operativo difícil de reconstruir a partir de los artículos: qué métrica da la alarma primero, qué tipo de loss spike es recuperable, cómo particionar los datos, cuándo reiniciar y qué optimizaciones solo funcionan en ciertas topologías. El razonamiento de Lample al emprender fue que el equipo ya había pagado meses de costo de trial-and-error en grandes empresas, por lo que podía lograr entrenamientos confiables más rápido en el entorno de una empresa emergente, con recursos más limitados.

\lecturefigure{08-data-heavy-team.png}{El equipo inicial de Mistral dedicó la mayor parte de su personal a los datos, no solo a escribir código de entrenamiento.}{Subtítulos locales 17:10--18:10; «1 persona en entrenamiento, 6 en datos» es la proporción organizacional mencionada de viva voz en clase.}

\begin{knowledgebox}{Lectura de la figura: por qué el data work se subestima con facilidad}
El trabajo con datos incluye auditoría de fuentes, deduplicación, filtrado por calidad, proporción de idiomas, detección de contaminación, análisis de ejemplos difíciles y diseño de evaluation slices. No tiene un nombre llamativo de nueva architecture, pero determina directamente la señal de entrenamiento. El docente incluso usa «alienating» para describir lo poco vistoso de este trabajo, pero enseguida enfatiza que es lo más crucial; esta es precisamente la teacher voice que debe conservarse.
\end{knowledgebox}

\teachervoice{En clase también se comparó la transferibilidad de la infraestructura de distintas grandes empresas: las plataformas internas de más alto nivel facilitan el trabajo diario, pero pueden ocultar los detalles de la planificación de tareas, el almacenamiento y la pila de entrenamiento; los entornos de más bajo nivel exigen que los investigadores resuelvan más problemas por sí mismos, lo que, al dejar la organización, facilita reconstruir el sistema. No se trata de alentar a reinventar la rueda, sino de recordar al equipo que conserve la explicabilidad y la capacidad operativa de las capas clave.}

\subsection{Resumen de la sección}

La dificultad del entrenamiento a gran escala proviene de nuevos regímenes numéricos, ciclos de retroalimentación costosos y la naturaleza tácita del conocimiento. Un equipo confiable no solo tiene más GPU: también necesita canary, observabilidad, registro de versiones, capacidad de trabajo con datos y la experiencia de ensayo y error que ya pagó.
